%%%%%%%%%%%%%%%%%%%%%%%%%%%%%%%%%%%%%%%%%%%%%%%%%%%%%%%%%%%%%
\section{Introduction}\label{sec:intro}
%%%%%%%%%%%%%%%%%%%%%%%%%%%%%%%%%%%%%%%%%%%%%%%%%%%%%%%%%%%%%
Benchmark problems are central to the development of fluid--solid interaction (FSI) codes: they are used to demonstrate accuracy, to compare coupling algorithms and to support claims that a new method is correct.
Many such problems exist, from the flag-behind-a-cylinder family of \citet{Turek2006} to pressure-pulse tubes, lid-driven cavities with flexible membranes and three-dimensional flows past elastic plates.
Far fewer of them come with a benchmark \emph{solution} whose own numerical error has been established: published values are often single-mesh results, are read from figures, are reported without time-step studies, or differ between sources by more than the discretisation errors that they are used to judge.
Agreement with such a value is then weak evidence, and disagreement is ambiguous: it may mean that the code under test is wrong, that the reference is wrong, or that the code's own solution is not yet resolved.
This paper addresses that gap for continuum FSI with an open, executable verification suite in which the quality of every reference and the spatial, temporal and iterative errors of every result are stated explicitly.

FSI methods can be classified along three axes: the modelling viewpoint (Eulerian, Lagrangian, or mixed), the degree of interface coupling (weak vs.\ strong), and the solver structure (partitioned vs.\ monolithic).

Eulerian approaches treat both solid and fluid domains as ``fluid-like'' (e.g.\ volume-of-fluid~\citep{jain2019conservative}, level-set~\citep{dunne2006eulerian, cottet2008eulerian}), easily allowing arbitrarily large solid deformations, but conservative interface descriptions and the inclusion of advanced solid constitutive laws remain challenging.
In contrast, Lagrangian (particle-type) approaches treat both domains as ``solid-like'', e.g.\ smooth-particle hydrodynamics, the material point method, and the particle finite element method (PFEM)~\citep{aubry2005particle, aubry2006fractional, cremonesi2020state}.
These methods are promising for large-deformation problems but typically come with a prohibitive computational cost.
Mixed Eulerian--Lagrangian approaches treat the solid (and interface) as Lagrangian and the fluid using an arbitrary Lagrangian--Eulerian (ALE) description.
ALE is effective for small to moderate interface deformations but eventually fails for very large deformations~\citep{wall2006large}, motivating methods that allow the solid to ``float over'' an Eulerian fluid grid, such as immersed boundary/domain~\citep{peskin1977numerical, peskin2002immersed, iaccarino2003immersed}, Chimera/overset~\citep{ge2003three, miller2014overset, li2021fluid, cane2021patient}, and cut-cell methods~\citep{pasquariello2016cut, tao2020new, xie2020cartesian}.
Of these, overset methods offer large-deformation freedom while retaining orthogonal interface cells suitable for resolving fluid boundary layers, albeit with the downside of greater implementation complexity.

In weakly coupled (loosely coupled) approaches, the fluid and solid are solved once per time step without iteration over the interface conditions~\citep{farhat2000two}.
This is computationally inexpensive but unstable for strongly coupled systems with significant added-mass effects~\citep{Causin2005, Forster2006}, such as incompressible flows interacting with light, flexible structures.
Strongly coupled approaches strictly enforce the interface kinematic and dynamic conditions at each time step; the present work focuses on the strongly coupled case.

Strong coupling can be enforced through either a partitioned or a monolithic solver structure.
In partitioned (staggered, two-system) approaches, separate solvers --- and potentially separate discretisations --- are used for the fluid and solid~\citep{Tukovic2018}, and outer iterations are performed over the interface at each time step.
These are accelerated using Aitken relaxation~\citep{Degroote2009, bungartz2016precice, irons1969version, kuttler2008fixed}, quasi-Newton schemes~\citep{santiago2020hpc, delaisse2022surrogate}, or Robin-type interface conditions~\citep{Tukovic2018b}.
The principal limitations are robustness and efficiency: for strongly coupled, highly nonlinear problems, outer-iteration convergence can be slow or fail entirely, particularly under strong added-mass loading.
In contrast, monolithic approaches treat the fluid and solid as a single system, solved concurrently within the same linearised system~\citep{Degroote2009}.
This avoids outer iterations and can handle tightly coupled nonlinear problems robustly, at the cost of reduced modularity and a poorly conditioned linear system that is more difficult to solve.

The verification literature for FSI is smaller than that for its component fields.
The proposal of \citet{Turek2006} defines the most widely used problems and reports multi-level reference values, and goal-oriented error estimation has been developed for FSI in the finite element setting~\citep{Richter2012, Richter2017}; monolithic finite element and hybrid finite volume/finite element codes have been demonstrated on these and related problems~\citep{Heil2004, Heil2008, Eken2016, Lozovskiy2019}, with convergence evidence that ranges from multi-level studies to single-mesh results.
Partitioned finite volume FSI, including the OpenFOAM-based framework of \citet{Tukovic2018} on which the solids4foam toolbox~\citep{Cardiff2025} builds, has mostly been assessed by comparison with such published values.
Across this literature, the error of the reference itself is rarely quantified, iterative (coupling) errors are seldom separated from discretisation errors, and the time integration of a reference is often first order with a large step, so that two codes can disagree because their references differ rather than because either is wrong.
General verification practice --- observed orders from systematic refinement, separation of error sources, and the distinction between verification and validation~\citetodo{Roache 1998; Oberkampf \& Roy 2010} --- is well established for single-field problems but is applied unevenly to coupled FSI.

The present work applies that practice to a suite of continuum FSI benchmarks, all solved with one partitioned, cell-centred finite volume framework: a PIMPLE-type ALE fluid, small-strain and total Lagrangian finite volume solids with standard and high-order discretisations~\citep{Cardiff2026, Batistic2026}, and Dirichlet--Neumann (Aitken~\citep{kuttler2008fixed}, IQN-ILS~\citep{Degroote2009}) or Robin--Neumann~\citep{Tukovic2018b} coupling.
The aim is not to promote a solver but to establish, case by case, what the available references can and cannot verify, and how much of the remaining difference is attributable to the code under test.
Verification here means convergence of the discrete solution to the solution of the stated mathematical problem; agreement with experimental data is validation, is kept separate, and is never used as evidence of convergence.

The specific contributions are:
\begin{enumerate}
    \item A reference-quality classification for FSI benchmarks (category A: analytical; B: independently converged numerical solution, possibly of a problem with a small, quantified model-form difference; C: published values with incompletely quantified error; C-weak: digitised, single-mesh, inconsistent or same-lineage values), applied to every case in the suite.
    \item Spatial, temporal and iterative error studies for each benchmark, with observed orders reported only where the evidence supports them, and with the slowest-converging quantity of interest identified.
    \item Evidence that verification exposes real defects and misleading benchmark agreement: two solver defects found through a cross-platform discrepancy in a three-dimensional tube, one compiler-dependent and one that changed the interface interpolation with mesh level; a two-level agreement with the Hron--Turek FSI3 reference that disappears on a third level; a published code-to-code discrepancy that is largely reproduced by the difference in time integration; and a default coupling convergence test that can accept unconverged steps.
    \item An open, executable implementation of the suite in solids4foam~\citep{Cardiff2025}: case set-ups, verification drivers that compute the quantities of interest and observed orders, and machine-readable reference files with their provenance, so that the studies can be re-run on other codes and platforms.
\end{enumerate}

Section~\ref{sec:math_model} states the governing equations and Section~\ref{sec:numerical_methods} the partitioned finite volume procedures used for every result.
Section~\ref{sec:verification_methodology} sets out the verification methodology, Section~\ref{sec:benchmark_suite} presents the benchmark suite, Section~\ref{sec:cross_benchmark} draws cross-benchmark conclusions, and Section~\ref{sec:recommended_suite} proposes a recommended verification suite.
Section~\ref{sec:discussion} discusses the findings and Section~\ref{sec:conclusion} concludes.
\hl{Before submission, verify all references have correct details}
